% ===== Document Preamble =====
\documentclass[11pt]{article}
\usepackage[margin=0.5in]{geometry}
\usepackage{amsmath,amssymb}
\usepackage{xcolor}
\usepackage{enumitem}
\usepackage{needspace}
\usepackage{comment}

% ===== Colors =====
\definecolor{AnsRed}{RGB}{170,0,0}
\definecolor{WrongBlue}{RGB}{0,0,160}
\definecolor{SkipGray}{RGB}{120,120,120}

% ===== Choice list styling (tight, consistent) =====
\setlist[enumerate]{%
  leftmargin=1.5em,
  itemsep=0.25em,
  topsep=0.25em,
  label=\textbf{\Alph*.}
}

% ===== Style Wrappers =====
\newcommand{\QuestionStyle}[1]{{\large\bfseries\boldmath #1}}
\newcommand{\AnswerStyle}[1]{{\large\bfseries\boldmath #1}}
\newcommand{\ReasonStyle}[1]{{\normalsize #1}}

% ===== Show / Hide Reasoning Toggle =====
\newif\ifshowreason
\showreasontrue      % Instructor version
% \showreasonfalse   % Student version

% ===== Choice Item Macros (Answer + Reasoning) =====
% Usage:
%   \Correct{Answer}{Reasoning}
%   \Wrong{Answer}{Reasoning}

\newcommand{\Correct}[2]{%
  \item \AnswerStyle{\textcolor{AnsRed}{#1}}%
  \ifshowreason\\[-0.2em]
    \ReasonStyle{\textcolor{AnsRed}{#2}}%
  \fi
}

\newcommand{\Wrong}[2]{%
  \item \AnswerStyle{\textcolor{WrongBlue}{#1}}%
  \ifshowreason\\[-0.2em]
    \ReasonStyle{\textcolor{WrongBlue}{#2}}%
  \fi
}

% ===== Multiple-Choice Question Block Macro =====
% Usage: \QBlock{Q1}{stem...}{<items>}{answer key}

\newcommand{\QBlock}[4]{%
  \Needspace{8\baselineskip}%
  \noindent
  \begin{minipage}{\linewidth}
    \QuestionStyle{#1.~#2}

    \vspace{0.25em}
    \begin{enumerate}
      #3
    \end{enumerate}

    \vspace{0.1em}
    \noindent\textcolor{AnsRed}{\textit{\textbf{Answer:}~#4}}

    \vspace{0.25em}
    \hrule height 0.4pt
  \end{minipage}
  \vspace{0.75em}
}

% ===== Open-Ended Question Block Macro =====
% Usage: \OEBlock{Q1}{stem...}{answer text}

\newcommand{\OEBlock}[3]{%
  \Needspace{5\baselineskip}%
  \noindent
  \begin{minipage}{\linewidth}
    \QuestionStyle{#1.~#2}

    \vspace{0.25em}
    \noindent\textcolor{AnsRed}{\textit{\textbf{Answer:}~#3}}

    \vspace{0.25em}
    \hrule height 0.4pt
  \end{minipage}
  \vspace{0.75em}
}

% ===== Multi-Part Question Block Macro =====
% Usage: \MPBlock{Q1}{stem...}{<parts>}

\newcommand{\MPBlock}[3]{%
  \Needspace{8\baselineskip}%
  \noindent
  \begin{minipage}{\linewidth}
    \QuestionStyle{#1.~#2}

    \vspace{0.25em}
    #3

    \vspace{0.25em}
    \hrule height 0.4pt
  \end{minipage}
  \vspace{0.75em}
}

% ===== MC Sub-Part Macro =====
% Usage: \MCPart{a}{stem...}{<items>}{answer key}

\newcommand{\MCPart}[4]{%
  \par\vspace{0.35em}
  \noindent\textbf{(#1)}~#2
  \begin{enumerate}
    #3
  \end{enumerate}
  \vspace{0.1em}
  \noindent\textit{\textbf{Answer:}~#4}
  \par
}

% ===== Open-Ended Sub-Part Macro =====
% Usage: \OEPart{a}{stem...}{answer text}

\newcommand{\OEPart}[3]{%
  \par\vspace{0.35em}
  \noindent\textbf{(#1)}~#2
  \par\vspace{0.15em}
  \noindent\textit{\textbf{Answer:}~#3}
  \par
}

% ===== Skipped Question Note Macro =====
% Usage: \SkipBlock{Q1}{reason}

\newcommand{\SkipBlock}[2]{%
  \noindent\textcolor{SkipGray}{\textit{[#1 omitted: #2]}}%
  \vspace{0.5em}\par
}

% ===== Document Start =====

\begin{document}

\begin{center}
{\huge \textbf{Physics Isomorphic Questions}}\\[0.25em]
{\Large Modeling, Units, Vectors, Kinematics, Forces, Momentum}\\[0.25em]
{\large Questions generated below. Multiple-choice answers in RED; incorrect answers in BLUE.}
\end{center}

\vspace{0.75em}

\QBlock{Q1}{Which of the following modeling statements is most appropriate?}{
  \Correct{We can safely treat a commuter train as a point mass when tracking its position along a 50 km rail line over time.}{When the length scale of interest (tens of kilometers) is much larger than the object's size (a train is a few hundred meters at most), its detailed shape is negligible for translational motion. In that case, using a point-mass model and tracking the center-of-mass position is appropriate.}
  \Wrong{We can safely treat the train as a point mass when analyzing the forces inside a coupler between two cars during acceleration.}{This neglects internal structure; analyzing coupler forces requires treating cars as extended bodies with contact forces, not a single point mass.}
  \Wrong{We can never assume a vacuum in mechanics problems, even when we want to ignore air resistance.}{This is a misconception: assuming vacuum (or ``negligible drag'') is a standard idealization used specifically to ignore air resistance.}
  \Wrong{None of the above statements is appropriate.}{Option A is a standard, justified approximation for large-scale position tracking, so ``none of the above'' is not correct.}
}{A}

\QBlock{Q2}{Which of the following unit-conversion equations is correct?}{
  \Wrong{$1\,\text{ms} = 1000\,\text{s}$}{This inverts the milli- prefix. The prefix ``milli-'' means $10^{-3}$, not $10^{3}$.}
  \Correct{$1\,\text{ms} = 0.001\,\text{s}$}{The prefix ``milli-'' means $10^{-3}$. Therefore $1\,\text{ms} = 10^{-3}\,\text{s} = 0.001\,\text{s}$.}
  \Wrong{$1\,\text{s} = 0.001\,\text{ms}$}{This reverses the conversion. Since $1\,\text{s}=1000\,\text{ms}$, it cannot equal $0.001\,\text{ms}$.}
  \Wrong{$1000\,\text{ms} = 0.001\,\text{s}$}{This applies the factor of 1000 in the wrong direction. In fact $1000\,\text{ms} = 1\,\text{s}$.}
}{B}

\QBlock{Q3}{Which option lists an incorrect SI unit for the quantity?}{
  \Wrong{Power: watt ($\text{W}$)}{Correct: power has SI units of watts, with $1\,\text{W}=1\,\text{J/s}$.}
  \Wrong{Momentum: $\mathrm{kg\,m/s}$}{Correct: momentum is $\vec p=m\vec v$, so its units are $\mathrm{kg\,m/s}$.}
  \Wrong{Electric charge: coulomb ($\text{C}$)}{Correct: the SI unit of charge is the coulomb.}
  \Correct{Pressure: $\mathrm{N\,m}$}{Pressure is force per area, so its SI unit is $\mathrm{N/m^2}$, also called the pascal (Pa). Writing $\mathrm{N\,m}$ (newton-meter) corresponds to torque/work units, not pressure.}
}{D}

\QBlock{Q4}{A rover is at the origin of an $x$-$y$ coordinate system. A beacon is located at position $\vec r = 8\hat\imath - 6\hat\jmath$ (units: kilometers). Which statement is correct?}{
  \Wrong{The distance from the rover to the beacon is $14\,\text{km}$.}{This resembles adding absolute component values ($8+6$). The distance is the magnitude $|\vec r|=\sqrt{8^2+(-6)^2}$, not $|8|+|{-6}|$.}
  \Wrong{The distance from the rover to the beacon is $8\,\text{km}$.}{This confuses the $x$-component with the full distance. The nonzero $y$-component also contributes via the Pythagorean theorem.}
  \Wrong{The distance from the rover to the beacon is $6\,\text{km}$.}{This confuses the magnitude with the $y$-component. The distance must include both components: $|\vec r|=\sqrt{x^2+y^2}$.}
  \Correct{None of the above is correct.}{The distance is $|\vec r|=\sqrt{8^2+(-6)^2}=\sqrt{64+36}=\sqrt{100}=10\,\text{km}$. Since $10\,\text{km}$ is not listed in A--C, the correct choice is none of the above.}
}{D}

\QBlock{Q5}{If the displacement of an object along the $x$-axis is given by
\[
\vec s(t)=\big((0.80\,\mathrm{m/s^3})t^3-(4.20\,\mathrm{m/s^2})t^2\big)\,\hat\imath,
\]
which statement is correct?}{
  \Wrong{This motion has constant acceleration.}{A cubic position function implies acceleration depends on time. Differentiating twice gives $a(t)=\dfrac{d^2s}{dt^2}=4.8t-8.4$, which is not constant.}
  \Wrong{The magnitude of the acceleration is $4.20\,\mathrm{m/s^2}$.}{This confuses a coefficient in $s(t)$ with acceleration. The acceleration comes from the second derivative, not directly from the $t^2$ coefficient.}
  \Correct{The magnitude of the acceleration at $t=0$ is $8.4\,\mathrm{m/s^2}$.}{Compute $v(t)=\dfrac{ds}{dt}=(2.40t^2-8.40t)\hat\imath$ and then $a(t)=\dfrac{dv}{dt}=(4.80t-8.40)\hat\imath$. At $t=0$, $a(0)=-8.40\,\mathrm{m/s^2}$ so the magnitude is $|a(0)|=8.4\,\mathrm{m/s^2}$.}
  \Wrong{The magnitude of the acceleration at $t=0$ is $0\,\mathrm{m/s^2}$.}{This typically comes from dropping the constant term after differentiating. Here $a(0)=-8.4\,\mathrm{m/s^2}$, not zero.}
}{C}

\QBlock{Q6}{A robot moves along a straight path from position $\vec r_A=-1.0\hat\imath+7.0\hat\jmath$ m to position $\vec r_B=5.0\hat\imath+2.0\hat\jmath$ m in $4.0\,\text{s}$. Which statement is correct?}{
  \Wrong{The robot travels a distance of $11\,\text{m}$.}{This is the mistake of adding component changes ($6+5$). The straight-line distance is the magnitude $\sqrt{(\Delta x)^2+(\Delta y)^2}$, not the sum.}
  \Wrong{The robot's speed at $t=0$ must be $1.95\,\text{m/s}$.}{Only knowing endpoints and total time does not uniquely determine the initial speed; many different time-dependent motions along the line can connect A to B in $4\,\text{s}$.}
  \Correct{The robot travels a distance of about $7.8\,\text{m}$.}{The displacement is $\Delta\vec r=(5-(-1))\hat\imath+(2-7)\hat\jmath=(6,-5)\,\text{m}$. The straight-line distance is $|\Delta\vec r|=\sqrt{6^2+(-5)^2}=\sqrt{61}\approx 7.81\,\text{m}\approx 7.8\,\text{m}$.}
  \Wrong{The acceleration at $t=0$ must be $0.49\,\text{m/s}^2$.}{This incorrectly assumes constant-acceleration kinematics without it being stated; acceleration cannot be determined from only the endpoints and elapsed time.}
}{C}

\QBlock{Q7}{A soccer ball is kicked with initial speed $v_0$ at an angle $\theta$ above the horizontal (ignore air resistance). What is the ball's speed at maximum height?}{
  \Correct{$v_0\cos\theta$}{At the top of the trajectory, the vertical velocity component is zero: $v_y=0$. The horizontal component remains constant without air resistance: $v_x=v_0\cos\theta$, so the speed is $v=\sqrt{v_x^2+v_y^2}=v_0\cos\theta$.}
  \Wrong{$v_0\sin\theta$}{This incorrectly uses the initial vertical component. At maximum height the vertical component has been reduced to $0$ by gravity.}
  \Wrong{$v_0$}{This assumes the speed stays equal to the launch speed. Gravity changes the vertical component, so the speed generally changes during flight.}
  \Wrong{$0$}{Only the vertical component becomes zero at the top; the ball still has horizontal speed $v_0\cos\theta$.}
}{A}

\QBlock{Q8}{A drone initially has velocity $\vec v_0=(15.0\,\mathrm{m/s})\hat\imath$ when its controller commands a constant acceleration $\vec a=(1.50\,\mathrm{m/s^2})\hat\imath+(3.00\,\mathrm{m/s^2})\hat\jmath$. Which statement is correct?}{
  \Wrong{The $x$-component of the drone's velocity remains constant.}{Since $a_x=1.50\,\mathrm{m/s^2}\neq 0$, $v_x$ changes with time via $v_x(t)=v_{0x}+a_x t$.}
  \Correct{Both the direction of motion and the speed change.}{With $\vec v(t)=\vec v_0+\vec a t$, the $y$-component grows from zero because $a_y\neq 0$, so the direction of $\vec v$ changes. Also, because $a_x\neq 0$, the speed $|\vec v|$ increases (it is not purely perpendicular acceleration).}
  \Wrong{The $y$-component of the drone's velocity remains constant.}{Because $a_y=3.00\,\mathrm{m/s^2}\neq 0$, the $y$-velocity evolves as $v_y(t)=v_{0y}+a_y t$ and cannot remain constant.}
  \Wrong{None of the above is correct.}{Option B correctly describes the consequences of a 2D acceleration with nonzero $a_x$ and $a_y$.}
}{B}

\QBlock{Q9}{Two point charges are separated by a distance $r$. If the magnitude of each charge is doubled and the separation distance is also doubled, what happens to the magnitude of the electrostatic force between them?}{
  \Wrong{It increases by a factor of $4$.}{Doubling both charges makes $|q_1q_2|$ four times larger, but doubling distance makes $r^2$ four times larger too; forgetting the $r^2$ factor causes this mistake.}
  \Wrong{It decreases by a factor of $4$.}{This accounts for the inverse-square dependence on $r$ but ignores that both charges were increased, which increases the numerator.}
  \Wrong{It increases by a factor of $2$.}{This is a common linear-scaling error (treating $F\propto q/r$ instead of $F\propto q_1q_2/r^2$).}
  \Correct{It stays the same.}{Coulomb's law gives $F=k\dfrac{|q_1q_2|}{r^2}$. Doubling both charges multiplies $|q_1q_2|$ by $4$, and doubling $r$ multiplies $r^2$ by $4$, so the factor is $4/4=1$ and the force is unchanged.}
  \Wrong{It decreases by a factor of $2$.}{This comes from mixing up linear and square dependence on $r$. The distance change affects $r^2$, not $r$.}
}{D}

\QBlock{Q10}{A planet's escape speed is given by $v_e=\sqrt{\dfrac{2GM}{R}}$, where $M$ and $R$ are the planet's mass and radius. Which statement is correct?}{
  \Wrong{A heavier spacecraft requires a higher escape speed from the same planet.}{This incorrectly inserts the spacecraft mass into the escape-speed requirement. The formula $v_e=\sqrt{2GM/R}$ contains only the planet's $M$ and $R$ (and $G$).}
  \Wrong{A lighter spacecraft requires a lower escape speed from the same planet.}{This repeats the misconception that the escaping object's mass matters. Escape speed is independent of the spacecraft mass in this model.}
  \Correct{Escape speed is higher for a planet with greater mass and smaller radius.}{From $v_e\propto \sqrt{M/R}$, increasing $M$ increases $v_e$ and decreasing $R$ also increases $v_e$. A more massive and more compact planet therefore has a larger escape speed.}
  \Wrong{Escape speed is higher for a planet with smaller mass and larger radius.}{This reverses the dependence: smaller $M$ and larger $R$ both decrease $M/R$, which lowers $v_e$.}
}{C}

\QBlock{Q11}{A constant force $\vec F=(5.60\,\mathrm{N})\hat\imath+(-2.40\,\mathrm{N})\hat\jmath$ acts on a particle as it moves along the $x$-axis from $x=-1.20\,\mathrm{m}$ to $x=3.30\,\mathrm{m}$. What is the work done by this force?}{
  \Correct{$25.2\,\mathrm{J}$}{The displacement is purely along $x$: $\Delta\vec r=(3.30-(-1.20))\hat\imath=4.50\hat\imath\,\mathrm{m}$. Work is the dot product $W=\vec F\cdot\Delta\vec r=F_x\Delta x=(5.60)(4.50)=25.2\,\mathrm{J}$; the $y$-component does no work because $\Delta y=0$.}
  \Wrong{$10.8\,\mathrm{J}$}{This is what you would get by (incorrectly) using $F_y\Delta x$ or mixing components. Since there is no $y$-displacement, the $F_y$ term cannot contribute.}
  \Wrong{$27.4\,\mathrm{J}$}{This likely comes from using the force magnitude $|\vec F|$ instead of the component parallel to motion. Work depends on $\vec F\cdot\Delta\vec r$, not $|\vec F||\Delta\vec r|$ unless the force is parallel.}
  \Wrong{$-25.2\,\mathrm{J}$}{The sign would be negative only if the component of force along the displacement were opposite the motion. Here $F_x>0$ and $\Delta x>0$, so $W=F_x\Delta x$ is positive.}
}{A}

\QBlock{Q12}{A $3.2\,\mathrm{kg}$ cart moving at $1.8\,\mathrm{m/s}$ collides with a bumper and rebounds with speed $0.9\,\mathrm{m/s}$ in the opposite direction. The collision lasts $0.15\,\mathrm{s}$. What is the magnitude of the average force on the cart during the collision?}{
  \Wrong{$38.4\,\mathrm{N}$}{This corresponds to using $\Delta v=1.8-0.9=0.9\,\mathrm{m/s}$ (subtracting speeds) instead of accounting for the velocity reversal, which makes $\Delta v$ larger.}
  \Correct{$57.6\,\mathrm{N}$}{Average force magnitude follows from impulse: $F_{\text{avg}}=\dfrac{|\Delta p|}{\Delta t}=\dfrac{m|\Delta v|}{\Delta t}$. With $v_i=+1.8\,\mathrm{m/s}$ and $v_f=-0.9\,\mathrm{m/s}$, $|\Delta v|=|v_f-v_i|=|-0.9-1.8|=2.7\,\mathrm{m/s}$, so $F_{\text{avg}}=\dfrac{(3.2)(2.7)}{0.15}=57.6\,\mathrm{N}$.}
  \Wrong{$19.2\,\mathrm{N}$}{This commonly comes from forgetting the division by collision time or using only one of the velocities. Impulse depends on the full momentum change over the stated time.}
  \Wrong{None of the above is correct.}{Option B correctly applies $F_{\text{avg}}=m|\Delta v|/\Delta t$ with the direction change included, so ``none'' is not correct.}
}{B}

\QBlock{Q13}{An $85\,\mathrm{kg}$ baseball catcher is initially at rest. A $145\,\mathrm{g}$ baseball traveling at $130\,\mathrm{km/h}$ is caught and sticks in the glove (perfectly inelastic). What is the catcher's speed immediately after the catch?}{
  \Correct{$0.061\,\mathrm{m/s}$}{Conserve momentum for a perfectly inelastic catch: $m v=(M+m)v_f$, so $v_f=\dfrac{m v}{M+m}$. Convert $130\,\mathrm{km/h}$ to $v=36.11\,\mathrm{m/s}$, then $v_f=\dfrac{(0.145)(36.11)}{85.145}\approx 6.15\times 10^{-2}\,\mathrm{m/s}\approx 0.061\,\mathrm{m/s}$.}
  \Wrong{$0.145\,\mathrm{m/s}$}{This resembles confusing the ball's mass in kilograms with the final speed, or neglecting the large catcher mass in the denominator. Because $M\gg m$, the final speed must be much smaller than $1\,\mathrm{m/s}$.}
  \Wrong{$36.1\,\mathrm{m/s}$}{This incorrectly assumes the catcher ends up moving at the ball's initial speed. Momentum conservation with $M\gg m$ implies $v_f\approx (m/M)v$, which is tiny.}
  \Wrong{$21.7\,\mathrm{m/s}$}{This comes from an incorrect mass ratio or unit conversion. A massive catcher cannot reach tens of m/s from catching a light baseball.}
}{A}

\end{document}
